\documentclass[11pt]{article}
\usepackage{coursenotes}
\begin{document}
\notesheader{7}{Observational Data}{When the logs must answer: adjustment, propensity scores, and double robustness}

\begin{decision}
No experiment was run. Treatment was assigned by a business process (a targeting rule, an account manager, an automated
system), and the logs record who got it and what happened. Can the logs say what the treatment did? Only if everything
that drove the assignment and also affects the outcome was recorded, and only for the kinds of unit that were sometimes
treated and sometimes not.
\end{decision}

\begin{goals}
  \item state the three identifying assumptions and read the assignment rule as the assignment mechanism;
  \item identify the ATE and ATT by standardization and by inverse-probability weighting, with proofs;
  \item show why a scalar propensity score suffices, and why predicting treatment well is a warning;
  \item choose controls with the backdoor logic, recognizing mediators, colliders, and M-bias;
  \item run the outcome-free design stage, diagnose overlap, and say how trimming changes the estimand;
  \item use the doubly robust (AIPW) estimator and say what double robustness does and does not protect against.
\end{goals}

%=====================================================================================
\sectionone

\subsection{Setup and estimands}

Treatment $D$, outcome $Y$, pre-treatment covariates $X$; $\mu_d(x) = \E[Y \mid D=d, X=x]$ and $e(x) = \Pr(D=1 \mid X=x)$. When a
recorded rule assigned treatment, $e$ is that rule viewed as a function of its inputs. This includes rules written by an
automated system: its decisions are the assignment mechanism, and the data it leaves are observational however sophisticated
it was. The ATE, ATT, and ATU differ whenever effects vary with $X$ and the rule's targeting depends on $X$; which one is
needed is a question about the decision.

\subsection{Identification}

\begin{assum}{Conditional exchangeability}{exch}
$\{Y(0), Y(1)\} \indep D \mid X$.
\end{assum}
\begin{assum}{Overlap}{overlap}
$0 < e(x) < 1$ for every $x$ in the population the estimand refers to.
\end{assum}
\begin{assum}{Consistency}{cons}
$Y = Y(D)$: one version of treatment, no interference.
\end{assum}

Exchangeability cannot be tested from the data it concerns. When a rule assigned treatment it has a concrete reading: the
rule's inputs are all in $X$, and whatever the rule left to chance was unrelated to outcomes. Defend it by finding out how
decisions were actually made, not by inspecting the data. Overlap is checkable and is where a decisive rule fails. For the ATT a weaker
pair suffices: $Y(0) \indep D \mid X$ and $e(x) < 1$.

\subsection{Standardization and weighting}

\begin{prop}{Standardization (the g-formula)}{std}
Under Assumptions~\ref{asm:exch}--\ref{asm:cons}, $\E[Y(d)] = \E_X[\mu_d(X)]$, so $\ATE = \E_X[\mu_1(X) - \mu_0(X)]$ and
$\ATT = \E[Y - \mu_0(X) \mid D = 1]$.
\end{prop}
\begin{proof}
$\E[Y(d)] = \E_X[\E[Y(d) \mid X]]$; exchangeability gives $\E[Y(d) \mid X] = \E[Y(d) \mid D=d, X]$ (defined by overlap), and
consistency replaces $Y(d)$ by $Y$. For the ATT, average over the treated, where $Y = Y(1)$.
\end{proof}

The ATT form needs only $\mu_0$, evaluated at the treated units' $x$, which is exactly where untreated units may be scarce.
With discrete $X$ the estimator is a weighted average of within-cell differences; with many covariates $\mu_d$ must be
modeled (``g-computation''), and the model will not report where it is extrapolating.

\begin{prop}{The IPW identity}{ipw}
Under Assumptions~\ref{asm:exch}--\ref{asm:cons}, $\E[DY/e(X)] = \E[Y(1)]$ and $\E[(1-D)Y/(1-e(X))] = \E[Y(0)]$. Under the
weaker ATT conditions,
\[
  \ATT = \frac{\E[DY] - \E\big[\tfrac{e(X)}{1-e(X)}(1-D)Y\big]}{\Pr(D=1)} .
\]
\end{prop}
\begin{proof}
$\E[DY(1) \mid X] = e(X)\E[Y(1) \mid X]$ by exchangeability; divide and average. For the ATT,
$\E[\tfrac{e}{1-e}(1-D)Y \mid X] = e(X)\E[Y(0) \mid X] = \E[DY(0) \mid X]$, so the numerator is $\E[D\{Y(1) - Y(0)\}]$.
\end{proof}

Weighting creates a pseudo-population in which $X$ no longer predicts $D$. Normalizing weights within each arm (the
\emph{H\'ajek} form) is preferred: it is invariant to shifting $Y$ and usually less variable. For the ATT, treated units keep
weight one and untreated units are weighted by the odds of treatment. The \emph{effective sample size} of an arm's weights,
$(\sum_i w_i)^2/\sum_i w_i^2$, says how many units the comparison really rests on.

\subsection{Why one number can replace all of \texorpdfstring{$X$}{X}}

\begin{prop}{The balancing property (Rosenbaum and Rubin)}{balance}
$D \indep X \mid e(X)$; hence if exchangeability holds given $X$, it holds given $e(X)$ alone.
\end{prop}
\begin{proof}
$\Pr(D=1 \mid X, e(X)) = e(X)$ and $\Pr(D=1 \mid e(X)) = \E[e(X) \mid e(X)] = e(X)$; they agree, so $D \indep X \mid e(X)$. Similarly
$\Pr(D=1 \mid Y(0), Y(1), e(X)) = \E[\Pr(D=1 \mid Y(0), Y(1), X) \mid Y(0), Y(1), e(X)] = e(X)$, using exchangeability given $X$.
\end{proof}

Adjusting for the scalar $e(X)$, by weighting, stratifying, or matching, does the work of the full vector. The propensity model
is not a prediction task: its job is to reproduce the assignment mechanism and balance $X$. A model that predicts $D$ almost
perfectly is reporting that the mechanism was nearly deterministic, which is a failure of overlap, not a success.

\subsection{What regression adjustment estimates}

$Y = \alpha + \tau D + X^\top\beta + \varepsilon$ is standardization with two restrictions: $\mu_0$ linear in $X$, and a constant
effect. If the first fails, residual confounding is attributed to $D$; if the second fails, the coefficient is a weighted
average of effects that matches none of the three estimands.

\begin{prop}{Variance-weighted estimand}{olsw}
If $\mu_0(X) = X^\top\beta$ exactly and effects vary, the OLS coefficient on $D$ converges to
$\E[\tau(X)\,e(X)(1-e(X))]/\E[e(X)(1-e(X))]$.
\end{prop}
\begin{proof}
Sketch, for $X$ saturated (so $e(X)$ is linear in $X$). Write $Y = X^\top\beta + D\tau(X) + \varepsilon$. By
Frisch--Waugh--Lovell the coefficient is $\E[\tilde D Y]/\E[\tilde D^2]$ with $\tilde D = D - e(X)$, which is orthogonal to $X^\top\beta$.
Given $X$, $\E[\tilde D D\tau(X) \mid X] = e(X)(1-e(X))\tau(X)$ and $\E[\tilde D^2 \mid X] = e(X)(1-e(X))$; average over $X$. Meeting~8 gives
the general version.
\end{proof}

The weights are largest where the rule was least decided ($e$ near $0.5$). Interacting $D$ with centered $X$, or standardizing
with a flexible model, targets the ATE or ATT instead.

\subsection{Good and bad controls}

Exchangeability is a statement about a particular $X$, and enlarging $X$ can destroy it (Meeting~1's backdoor criterion).
\begin{itemize}
  \item \textbf{Confounders} affect both $D$ and $Y$; a recorded rule's inputs are confounders by construction.
  \item \textbf{Mediators} lie on $D \to M \to Y$; adjusting removes the part of the effect through $M$.
  \item \textbf{Colliders} are caused by $D$ (or a cause of $D$) and by a cause of $Y$; adjusting creates an association
  (Meeting~1, collider result).
  \item \textbf{Pure predictors of treatment} (instruments) add variance and can amplify the bias from any remaining
  confounder.
\end{itemize}
Measure controls before the treatment decision. Timing is necessary but not sufficient: a pre-treatment variable can be a
collider if two unrecorded causes, one of $D$ and one of $Y$, both drive it (\emph{M-bias}). Justify each control by its causes,
not only its date. ``Everything in the warehouse'' almost always contains post-treatment variables.

\subsection{Overlap, trimming, and the estimand}

Diagnose overlap with the distribution of $\hat e(X)$ in each arm (regions where one arm is empty have no comparison) and with
the largest weights. \emph{Trimming} restricts to $\hat e(X) \in [a, 1-a]$, often with $a = 0.05$ or $0.10$. It stabilizes the estimate
but \emph{changes the estimand} to the average effect in the trimmed population, which must be reported as such (``accounts with fewer
than ten users, all of which were called, are excluded''). Standardization does not fail visibly in the same region: the outcome
model extrapolates silently. Compute both; where they disagree, look first at the overlap diagnostics.

\subsection{Design before analysis}

Every choice so far can be made without looking at an outcome, and should be (Rubin, 2008).
\begin{enumerate}
  \item \textbf{Reconstruct the assignment mechanism}: estimate $e(x)$ from the rule's inputs, or write it down if the rule is
  known.
  \item \textbf{Check overlap} in each arm and fix the trimming rule.
  \item \textbf{Fix the estimand} the trimmed population implies, and the estimator.
  \item \textbf{Check balance} of $X$ after weighting or matching (standardized mean differences).
  \item Record the design in writing, with a date. \textbf{Only then open the outcomes.}
\end{enumerate}
This is the observational counterpart of the frozen partition (Meeting~4) and the frozen list (Meeting~5).

\subsection{Doubly robust estimation}

\begin{prop}{AIPW and double robustness}{aipw}
Let
\[
  \phi = \tilde\mu_1(X) - \tilde\mu_0(X) + \frac{D\,(Y - \tilde\mu_1(X))}{\tilde e(X)} - \frac{(1-D)(Y - \tilde\mu_0(X))}{1 - \tilde e(X)}
\]
for working models $\tilde\mu_d$ and $\tilde e$. Under Assumptions~\ref{asm:exch}--\ref{asm:cons}, $\E[\phi] = \ATE$ if \emph{either}
$\tilde\mu_d = \mu_d$ for both $d$ \emph{or} $\tilde e = e$.
\end{prop}
\begin{proof}
Given $X$, $\E[\phi \mid X] - \tau(X) = (\mu_1 - \tilde\mu_1)(e/\tilde e - 1) - (\mu_0 - \tilde\mu_0)\{(1-e)/(1-\tilde e) - 1\}$; each term is a
product that vanishes if the outcome model or the propensity is correct.
\end{proof}

The first two terms are standardization; the weighted residuals correct what the outcome model got wrong. When both nuisances
are estimated flexibly, the product structure makes the estimator's error depend on the product of the two models' errors, and
with cross-fitting (Meeting~8) it has a valid normal approximation where neither IPW nor standardization alone does. Double
robustness protects against \emph{modeling} errors, not against unmeasured confounding: if exchangeability fails, every
estimator in these notes is biased.

\paragraph{Inference.} Bootstrap over units, re-estimating $\hat e$ (and $\hat\mu_d$) in each replicate, with clusters resampled as
blocks; or use the cross-fitted AIPW standard error of Meeting~8.

\subsection{How well does adjustment work? The LaLonde lesson}

LaLonde (1986) replaced a job-training experiment's control group with survey respondents and asked whether adjustment recovered
the experimental answer. The regressions of the time missed badly, sometimes with the wrong sign. Later work found propensity
methods do well when the comparison group overlaps the treated \emph{and} the covariates include prior values of the outcome,
and badly otherwise. The covariates matter more than the estimator.

\subsection{A first look at sensitivity}

If an unmeasured $U$ shifts the outcome by $\gamma$ per unit and differs between treated and untreated units with the same $X$ by
$\delta$, a linear adjustment is off by about $\gamma\delta$. Ask the reverse question: how large must $\gamma\delta$ be to move the
estimate below the break-even? If that needs a confounder stronger than the strongest one measured, the conclusion is robust.
Meeting~8 makes this precise.

\begin{pitfall}[adjustment errors]
Adjusting for post-treatment or collider variables; trimming without restating the estimand; choosing the specification after
seeing the outcome; treating a near-perfect propensity model as good news; believing double robustness covers unmeasured
confounders.
\end{pitfall}

%=====================================================================================
\sectiontwo

\paragraph{Setting.} FitLife's membership staff decide which members to call before their renewal date. A call costs ¥30 of staff
time; a renewal is worth ¥500, so the break-even is $0.06$. Staff mostly call members with high recent attendance. Attendance in
the prior quarter, recorded before the call list is made, is the only input staff use (an assumption checked by interviewing them).
\begin{center}
\begin{tabular}{lcccccc}
\toprule
Attendance & Members & Called & $\hat e$ & Renewal if called & Renewal if not & Effect \\
\midrule
High & 1{,}000 & 800 & 0.8 & 0.90 & 0.85 & 0.05 \\
Low  & 1{,}000 & 100 & 0.1 & 0.60 & 0.50 & 0.10 \\
\bottomrule
\end{tabular}
\end{center}

\paragraph{Step 1: design, before outcomes.} Both strata have called and uncalled members, so overlap holds; no trimming. ATE
weights give an effective sample of $356$ of the 900 called and $655$ of the 1{,}100 uncalled; ATT weights on the uncalled
($e/(1-e) = 4$ for high, $1/9$ for low) give $252$ of 1{,}100. The ATT rests mostly on the 200 uncalled high-attendance members.

\paragraph{Step 2: naive.} Called members renew at $(800 \times 0.90 + 100 \times 0.60)/900 = 0.867$, uncalled at
$(200 \times 0.85 + 900 \times 0.50)/1{,}100 = 0.564$: a difference of $0.303$, five times the break-even. It is the ATT ($0.056$) plus
selection bias ($0.247$).

\paragraph{Step 3: standardization.} $\ATE = 0.5 \times 0.05 + 0.5 \times 0.10 = 0.075$;
$\ATT = (800 \times 0.05 + 100 \times 0.10)/900 = 0.056$; $\ATU = (200 \times 0.05 + 900 \times 0.10)/1{,}100 = 0.091$.

\paragraph{Step 4: IPW.} Weights $1/0.8 = 1.25$ (called, high), $1/0.1 = 10$ (called, low), $1/0.2 = 5$ (uncalled, high), $1/0.9$ (uncalled,
low):
\[
  \widehat{\E[Y(1)]} = \frac{800 \times 1.25 \times 0.90 + 100 \times 10 \times 0.60}{2{,}000} = 0.750,
\]
\[
  \widehat{\E[Y(0)]} = \frac{200 \times 5 \times 0.85 + 900 \times \tfrac{1}{0.9} \times 0.50}{2{,}000} = 0.675 .
\]
IPW ATE $= 0.075$, as it must be with a saturated propensity model.

\paragraph{Step 5: regression.} OLS of renewal on a call dummy and a high-attendance dummy weights strata by
$n_s\,e_s(1-e_s)$: $160$ for high, $90$ for low. So $\hat\tau_{\text{OLS}} = (160 \times 0.05 + 90 \times 0.10)/250 = 0.068$, neither the ATE nor
the ATT (Result~\ref{prop:olsw}).

\paragraph{Step 6: double robustness.} Suppose the analyst's outcome model ignores attendance and predicts the pooled rates
($0.867$ called, $0.564$ uncalled), but the propensity is right. In each stratum the AIPW correction terms have means
$\mu_{1s} - 0.867$ and $\mu_{0s} - 0.564$, which exactly undo the outcome model's error, so AIPW returns $0.075$. With the right outcome
model and a wrong propensity (say $0.45$ for everyone), the residuals have mean zero and AIPW again returns $0.075$.

\paragraph{The decision.} Calls as practiced earn an ATT of $0.056$, below the break-even of $0.06$: they lose money. Calling
low-attendance members (effect $0.10$) would pay; calling high-attendance ones ($0.05$) does not. The staff habit that inflated the
naive number is also what makes the program unprofitable.

\begin{exercise}[computation]
FitLife adds a third stratum: 500 ``elite'' members, all of whom were called, who renew at $0.97$. (a) What happens to overlap and to
the ATE? (b) What estimand can still be reported, and what is it?
\end{exercise}
\answer{(a) $e = 1$ for elite members: no uncalled elite member exists, so the population ATE (and the ATT, which includes elite
members) is not identified; IPW weights would be undefined. (b) Report the effect for the overlap population (high and low
attendance): ATE $0.075$ and ATT $0.056$ as before, stating that elite members are excluded and that their effect is unknown.}

\begin{exercise}[judgment]
An analyst proposes adding ``number of classes booked in the two weeks after the call date'' as a control, since it ``predicts
renewal very well''. What kind of variable is it, and what happens to the estimate?
\end{exercise}
\answer{It is measured after treatment and is likely affected by the call (calls prompt bookings): a mediator, and possibly a collider
if unmeasured motivation drives both bookings and renewal. Adjusting removes the part of the effect that works through bookings and
can open a non-causal path, biasing the estimate (usually toward zero). Strong prediction of the outcome is not a reason to adjust.}

%=====================================================================================
\sectionthree

\begin{extension}[Matching and weighting refinements]
Matching on $\hat e$ or on $X$ targets the ATT directly and discards unmatched controls. Overlap weights ($1 - e$ for treated, $e$ for
controls) target the overlap population and never explode. Entropy balancing and covariate-balancing propensity scores choose weights
that balance covariate moments exactly.
\end{extension}

\begin{extension}[Continuous dose]
With a continuous treatment the average dose--response $\mu(d) = \E_X[\E[Y \mid D=d, X]]$ is identified under the same three assumptions,
with overlap meaning that treatment varies within every region of $X$. If $\E[Y \mid D, X] = \theta D + g(X)$ the slope is the partially linear
model of Meeting~8.
\end{extension}

\subsection*{Annotated reading}
\begin{readinglist}
\reading{Rosenbaum and Rubin (1983), ``The central role of the propensity score in observational studies'',
\emph{Biometrika}}{The balancing property and its consequences}{Sections 1--3}{2}
\reading{Rubin (2008), ``For objective causal inference, design trumps analysis'', \emph{Annals of Applied
Statistics}}{The case for an outcome-free design stage}{All}{1}
\reading{Crump, Hotz, Imbens, and Mitnik (2009), ``Dealing with limited overlap'', \emph{Biometrika}}{How to choose a trimming
rule, and what estimand results}{Sections 1--3}{3}
\reading{LaLonde (1986), \emph{AER}; Dehejia and Wahba (1999), \emph{JASA}}{Can observational methods recover an experiment?}{Intro,
results}{1}
\reading{Imbens (2015), ``Matching methods in practice: three examples'', \emph{Journal of Human Resources}}{A worked tour of
design-stage diagnostics}{All}{2}
\reading{Handout H4, \emph{Instrumental variables}}{What to do when an unmeasured confounder makes exchangeability
implausible}{All}{2}
\end{readinglist}

\end{document}
